\chapter{Comparing options and ranking}\label{ch:options}

The \ui{Options} tab shows the schedules worth considering, ranked under the period's policy. Its job is to
make the decision quick: most of the time, one of the first two cards is the schedule you will commit.

\shot{options}{The Options tab of ``Weeks 48-49'' with two recommended options.}{options}

\begin{callouts}
  \item[1] \textbf{What matters most this period?} The ranking policy: \ui{Service first}, \ui{Balanced} or
           \ui{Cost first} (Section~\ref{sec:ranking}).
  \item[2] \textbf{How this policy scores schedules}: opens the exact scoring rule.
  \item[3] \textbf{Recommended option cards}, ranked \#1, \#2, \dots\ under the selected policy.
  \item[4] \textbf{What really differs}: one sentence on where the recommended options are practically equal
           and where they really differ --- often the fastest way to decide.
  \item[5] \textbf{Viewing bar}: the option currently selected, its key figures, and \ui{Commit option}. It
           stays at the bottom of the screen on every tab.
\end{callouts}

\section{Which schedules are recommended}\label{sec:recommended}

Every schedule is placed on three axes:

\begin{enumerate}
  \item \textbf{priority-weighted lateness} (minutes late $\times$ priority, summed over orders);
  \item \textbf{safety-stock shortfall} (units below safety, summed over end-of-day checkpoints);
  \item \textbf{operating cost} (the weighted sum of all other policy terms: shifts, setups, idle time,
        maintenance, surplus stock, \dots).
\end{enumerate}

A schedule is \emph{dominated} when another schedule is at least as good on all three axes and better on one.
PlanWise recommends only the schedules that are \textbf{not} dominated: every other schedule is beaten by a
recommended one on all three axes. Dominated schedules are not hidden --- they appear in the Analysis tab and
in the \ui{Compare with} list --- but there is no reason to commit one.

Recommended options are named after what they are best at:

\begin{center}
\begin{tabularx}{\linewidth}{>{\bfseries}p{4.2cm}X}
\toprule
\rowcolor{white}Name & Meaning \\
\midrule
Best on every axis & Only one schedule is recommended: it is at least as good as every other on all axes. \\
Protects priority orders & Lowest priority-weighted lateness. It may still have more raw hours late, but on low-priority orders. \\
Least shortfall & Lowest safety-stock shortfall. \\
Lowest cost & Lowest operating cost. \\
Middle ground & Recommended without being best on any single axis --- a compromise. \\
\bottomrule
\end{tabularx}
\end{center}

Names combine when an option wins several axes (``Least shortfall · Lowest cost''). Each card shows who built
it (\ui{Built by: Hybrid: search, then CP-SAT}, \ui{Manual edit}, \ui{Optimised around an edit}), the four
headline figures (green = best among the recommended options), a one-line comparison with \#1, and badges such
as \ui{Committed} and \ui{Viewing}. Click a card for a \ui{Quick look} at the schedule, or to make it the
viewed option.

If you edited schedules by hand and an edit is not recommended, it is listed under \ui{Your versions not
recommended} with the reason, so your work never disappears.

\section{Ranking policies}\label{sec:ranking}

The recommended options are ordered by a score. The switch at the top picks how the weighted terms of the
factory policy are combined:

\begin{description}[style=nextline, leftmargin=1.2em, font=\bfseries\color{brand}]
  \item[Service first] (default) Lateness and safety-stock shortfall are compared first (tier~1); operating
        cost only breaks ties (tier~2). A score reads ``tier~1 · tier~2''.
  \item[Balanced] Every weighted term is added into one score: trades a little safety stock or lateness for
        lower cost when the weights allow it.
  \item[Cost first] Operating cost is compared first; lateness and shortfall only break ties. Use it in slack
        weeks when orders can slip.
\end{description}

Switching re-ranks the same schedules instantly; nothing is recomputed. \ui{How this policy scores
schedules} shows the formula with the actual weights:

\shotf[0.82\linewidth]{options_scoring}{The scoring rule of ``Service first'' with the weights of the factory policy.}{options_scoring}

\begin{note}
The weights come from the factory's policy (Section~\ref{sec:policies}). This switch only changes the priority
order of the terms, not the weights. To change the weights, edit the policy on the Factory page or pick
another policy under Prepare.
\end{note}

\section{The trade-off chart}

Below the cards, a scatter plot places every schedule by operating cost (horizontal) and safety-stock
shortfall (vertical). Closer to the bottom-left corner is better. Named green dots are recommended options,
grey dots are other schedules, and a dashed line joins the recommended options.

\shotf[0.8\linewidth]{options_tradeoff}{Trade-off between operating cost and safety-stock shortfall.}{options_tradeoff}

\section{Key figures used everywhere}\label{sec:kpis}

The same figures appear on the cards, in the Schedule \& compare tiles, in the viewing bar and in exports:

\begin{longtable}{>{\bfseries}p{4.1cm}p{10.4cm}}
\toprule
\rowcolor{white}Figure & Definition \\
\midrule
\endhead
On time & Orders delivered by their due date, out of the orders of the period (for example 40/44). \\
Total lateness & Hours late summed over all late orders. \\
Priority-weighted lateness & Minutes late $\times$ order priority, summed (the policy's ``weighted tardiness''). \\
Safety-stock shortfall & Units below safety stock, summed over the end-of-day checkpoints. \\
Operating cost & Weighted sum of the non-service terms of the policy: non-productive time, setup, mold maintenance, makespan, surplus stock, pull-ahead, semi-finished left at the end. A relative figure for comparing schedules, not money. \\
Machine shifts open & Number of machine-shifts the schedule uses. Fewer shifts mean fewer people and less energy. \\
Idle time & Available minutes left unused inside opened shifts. \\
Setup & Changeover minutes on machines. \\
Mold maintenance & Mold service minutes. \\
Last required lot done & When the last lot needed by an order of the period leaves test and packing. \\
\bottomrule
\end{longtable}
